\chapter{Polar Curve}
\label{chap:polar-curve}
\lecture{3}{28 Apr. 22:52}{Polar Curves}

\begin{definition}[Polar Curve]
  \label{def:polar-curve}
  \begin{equation}
    \label{eqn:polar-curve}
    \begin{displaystyle}
      r = f(\theta)
    \end{displaystyle}
  \end{equation}
\end{definition}

\section{Symetricity of \(r=f(\theta)\)}
\label{sec:symetricity-of-r-f-theta}

\begin{enumerate}
  \item Given \((r,\theta)\) lies in \(r=f(\theta)\),
    \[
      (r,-\theta) \text{ or } (-r,\pi-\theta) \text{ in  } r
      = f(\theta) \implies r = f(\theta) \text{ is symmetric
      abt. x-axis. }
    \]
    \begin{eg}
      \(r=1+\cos\theta\) \quad (true)\\
      \(r=2a\sin\theta\) \quad (false)\\
      \(r=2a\sin2\theta\) \quad (false)
    \end{eg}

  \item Given \((r,\theta)\) lies in \(r=f(\theta)\),
    \[
      (-r,-\theta) \text{ or } (r,\pi-\theta) \text{ in  } r
      = f(\theta) \implies r = f(\theta) \text{ is symmetric
      abt. y-axis. }
    \]
    \begin{eg}
      \(r=1+\cos\theta\) \quad (false)\\
      \(r=2a\sin\theta\) \quad (true)\\
      \(r^{2}=\cos2\theta\) \quad (true)
    \end{eg}

  \item Given \((r,\theta)\) lies in \(r=f(\theta)\),
    \[
      (-r,\theta) \text{ or } (r,\pi+\theta) \text{ in  } r
      = f(\theta) \implies r = f(\theta) \text{ is symmetric
      abt. pole. }
    \]
    \begin{eg}
      \(r=1+\cos\theta\) \quad (false)\\
      \(r=\sin\theta\) \quad (false)\\
      \(r=\sin2\theta\) \quad (true)\\
      \(r^2=\sin2\theta\) \quad (true)
    \end{eg}

\end{enumerate}

\begin{table}[htpb]
  \centering
  \resizebox{0.8\textwidth}{!}{ 
    \begin{tabular}{|c|cc|}
      \hline
      \textbf{x-axis}&\((r,-\theta)\)&\((-r,\pi-\theta)\)\\
      \hline
      \textbf{y-axis}&\((r,\pi-\theta)\)&\((-r,-\theta)\)\\
      \hline
      \textbf{pole}&\((r,\pi+\theta)\)&\((-r,\theta)\)\\
      \hline
    \end{tabular}
  }
  \caption{Summary of Symmetricity of \(r=f(\theta)\)}
  \label{tab:summary-of-symmetricity-of-r-f-theta-}
\end{table}

\section{Slope of \(r=f(\theta)\)}
\label{sec:slope-of-r-f-theta}
Given a polar curve \( r = f(\theta) \), we can express the Cartesian coordinates as
\[
  x = r\cos\theta = f(\theta)\cos\theta, \qquad y = r\sin\theta = f(\theta)\sin\theta.
\]

The slope of the tangent to the curve at a point is
\[
  \frac{dy}{dx} = \frac{\frac{dy}{d\theta}}{\frac{dx}{d\theta}}.
\]

Differentiating \( x \) and \( y \) with respect to \( \theta \) (using the product rule):
\begin{align*}
  \frac{dx}{d\theta} &= \frac{d}{d\theta}\left( f(\theta)\cos\theta \right) = f'(\theta)\cos\theta - f(\theta)\sin\theta, \\
  \frac{dy}{d\theta} &= \frac{d}{d\theta}\left( f(\theta)\sin\theta \right) = f'(\theta)\sin\theta + f(\theta)\cos\theta.
\end{align*}

Therefore, the slope at \( \theta = \theta_0 \) is
\[
  \left. \frac{dy}{dx} \right|_{\theta = \theta_0}
    = 
    \left.
      \frac{f'(\theta)\sin\theta + f(\theta)\cos\theta}
      {f'(\theta)\cos\theta - f(\theta)\sin\theta}
      \right|_{\theta = \theta_0}.
    \]

    \bigskip

    \textbf{Slope at the Pole:}

    The \emph{pole} corresponds to \( r = f(\theta_0) = 0
    \). At the pole, the expressions simplify:
    \begin{align*}
      \left. \frac{dy}{dx} \right|_{r=0,\,\theta = \theta_0}
&= \frac{f'(\theta_0)\sin\theta_0}{f'(\theta_0)\cos\theta_0} \\
&= \tan\theta_0 \qquad \text{(provided \( f'(\theta_0) \neq 0 \)).}
    \end{align*}

    Thus, \textbf{the slope of the tangent at the pole is \(
    \tan\theta_0 \)}, provided the curve passes through the
    pole at \( \theta_0 \) and \( f'(\theta_0) \neq 0 \).

    \section{Graphing a Polar Curve}
    \label{sec:graphing-a-polar-curve}
    \begin{exercise}
      Plot the curve \(r=1+\cos\theta\), a cardiod(heart-shaped).
    \end{exercise}
    \begin{answer}
      The following is the process for graphing given curve.
      \begin{enumerate}
        \item Find the slope at pole:\\
          At pole, \(r=0\)\\
          We have,
          \begin{align*}
            r &= 1 + \cos\theta \\
            \text{Set } r = 0: \quad 0 &= 1 + \cos\theta \\
            \cos\theta &= -1 \\
            \theta &= \pi, \quad n \in \Z\\
            \therefore \tan\theta_o &= 0
          \end{align*}

        \item Find the Symmetricity of the curve:\\
          The curve is symmetric about only x-axis.

        \item Make \(r-\theta\) table for specific points.

      \end{enumerate}
    \end{answer}

    \begin{table}[htpb]
      \centering
      \begin{tabular}{|c|c|c|c|c|c|c|c|c|c|}
        \hline
        \(\theta\) & \(0\) & \(\frac{\pi}{6}\) &
        \(\frac{\pi}{4}\) & \(\frac{\pi}{3}\) &
        \(\frac{\pi}{2}\) & \(\frac{2\pi}{3}\) &
        \(\frac{3\pi}{4}\) & \(\frac{5\pi}{6}\) &
        \(\pi\)\\

        \hline

        \(r = 1+\cos(\theta)\) & \(2\) & \(1 +
        \frac{\sqrt{3}}{2}\) & \(1 +
        \frac{\sqrt{2}}{2}\) & \(1 + \frac{1}{2}\) &
        \(1\) & \(1 - \frac{1}{2}\) & \(1 -
        \frac{\sqrt{2}}{2}\) & \(1 -
        \frac{\sqrt{3}}{2}\) & \(0\)\\

        \hline
      \end{tabular}
      \caption{r-\(\theta\) Table}
      \label{tab:r-theta-table}
    \end{table}

    \textbf{Plot:}
    \begin{figure}[H]
      \centering
      \begin{polargrid}[2][2][12][2.2][1.2]
        \draw[domain=0:180,smooth,variable=\t,red,thick]
        plot ({(1+cos(\t)) * cos(\t)}, {(1+cos(\t)) * sin(\t)});
        \draw[domain=180:360,smooth,variable=\t,blue,thick]
        plot ({(1+cos(\t)) * cos(\t)}, {(1+cos(\t)) * sin(\t)});
      \end{polargrid}
      \caption{Plot of \(r=1+\cos\theta\)}
      \label{fig:plot-of-r-1-cos-theta}
    \end{figure}

